% =====================================================================
%  Chapter 12: Temperature and chemical potential
%  Gattringer & Lang, "Quantum Chromodynamics on the Lattice" (LNP 788)
%  Reconstructed from extract/ch12.txt; equations auto-numbered.
%  中文译本：第 12 章 温度与化学势
% =====================================================================

\chapter{温度与化学势}

自大爆炸以来，宇宙已经大幅冷却。尽管如此，从第一性原理理解高温和高物质密度下的 QCD 仍具有多方面的重要意义。其中一个方面显然是为了更好地理解大爆炸后不久夸克与胶子物质的行为，以及冷却过程中强子的凝聚。宇宙中的许多天体，如中子星，据推测具有足够高的密度或温度，因而可以预期强子物质的行为与普通原子有所不同。在地球上，高能重离子碰撞也提供了在极端条件下研究物质的可能性。因此，从第一性原理研究高温高密度下的 QCD 是格点方法的一项富有挑战性的任务。

直观上可以预期，在密度高于核子内部密度、温度高于介子质量时，强子作为束缚态将失去其个体性，并涌现出新的结构相。人们预期会发生相变，而格点表述正是研究这类集体现象的天然框架。这些研究应当揭示向不同物质态转变的存在与否及其类型、这些物质态的性质，以及可能存在的其他物质态在非常规环境中的行为，从而帮助我们理解实验以及天体物理观测。格点热力学进展的综述可见于文献 [1--5]。

在当今的格点方法中，我们总是研究处于平衡态的体系，而不是体系的演化和衰变动力学。这一挑战留待将来解决。

\section{温度的引入}

在第~1.4.3 节中，我们讨论了欧几里得量子场论与温度 $T$ 的热浴中经典统计力学体系的配分函数 (1.13) 之间的联系。对于量子力学情形，配分函数由下式给出
\begin{equation}
  Z(T) = \tr\!\left[e^{-H/(k_B T)}\right] = \tr\!\left[e^{-\beta H}\right],
\end{equation}
其中 $\hat{H}$ 是哈密顿算符，$\beta$ 现在表示逆温度 $\beta = 1/(k_B T)$，$k_B$ 为玻尔兹曼常数。这是标准的记号，但遗憾的是，它常常导致将 $\beta$ 与逆规范耦合混淆。除非另有说明，本章中 $\beta$ 专指逆温度。在随后的讨论中，我们采用通常的单位选择，令 $k_B = 1$。于是温度 $T$ 以能量或质量为单位给出，且 $\beta = 1/T$。

由于 (12.1) 中的迹，我们被限制在时间方向上周期（玻色子）或反周期（费米子）的场。与第~1 章类似，配分函数可以变换为对这类场位形的路径积分。然而，与那里的讨论相反，我们现在并不假设时间范围趋于无穷。相反，我们得到泛函积分（$\Phi$ 为某个一般的场）
\begin{equation}
  Z(T) = \int \mathcal{D}[\Phi]\,e^{-S_E[\Phi]},
\end{equation}
其中积分是对在有限时间方向上（反）周期的场进行的。欧氏作用量由对整个空间但有限时间范围的积分给出：
\begin{equation}
  S_E[\Phi] = \int_0^{\beta}\!\dd t \int_{\mathbb{R}^3}\!\dd^3 x\,
    \mathcal{L}_E\big(\Phi(t,x),\partial_\mu\Phi(t,x)\big).
\end{equation}

测度 $\mathcal{D}[\Phi]$ 和作用量 $S_E[\Phi]$ 照常以通常方式在格点上离散化。在离散化中我们发现，改变的只是我们的视角。到目前为止，在“零温”下工作时，我们感兴趣的是无限时空体积极限下的结果，并将偏差作为有限体积效应来讨论。空间和时间范围被认为远大于系统中最大的关联长度——即介子的关联长度。现在，空间仍被认为处于该极限下，但时间的物理范围被限制为 $\beta$。对于有限格点，空间范围为 $aN$，时间范围为 $aN_T$。因此我们有
\begin{equation}
  \beta = a\,N_T = \frac{1}{T}
\end{equation}
并且发现极限 $\beta\to\infty$ 对应于 $T\to 0$。我们将此解释为具有有限空间体积和固定温度 $T$ 的体系。这样一个体系的连续极限对应于在保持 $aN$ 和 $aN_T$ 固定的同时 $a\to 0$。当 $N/N_T$ 增大时，有限体积效应变小。

再次依据物理直觉，我们可以预期当介子的关联长度 $1/M_\pi$ 达到时间范围 $aN_T$ 时，将观察到有限温度效应；并且已发现 $T$ 确实与 $M_\pi$ 的数值相去不远。

时间范围 $\beta$ 保持有限这一事实具有重要的物理后果。如果将傅里叶变换应用于有限时间方向，则会发现只允许间隔为 $\Delta = 2\pi/\beta$ 的分立能级。该值 $\Delta$ 的整数倍通常被称为马祖巴频率 (Matsubara 频率)。由于格点结构，马祖巴频率不仅是分立的，而且被限制在 $(-\pi/a,\pi/a]$ 范围内，对玻色子态取值为 $k\Delta$（其中 $-N_T/2 + 1 \le k \le N_T/2$），$\Delta = 2\pi T = 2\pi/(aN_T)$。对于满足时间方向反周期边界条件的费米子，其取值平移 $\Delta/2$（参见附录~A），因此最小的夸克马祖巴模式能量为 $\Delta/2$。

\subsection{纯规范理论的分析}

让我们首先讨论有限温度下的纯规范理论。对于这种情况，我们在第~3.3.5 节中引入了一个重要的可观测量，即波利亚科夫圈 (Polyakov 圈)。它定义为时间方向上的规范链变量有序乘积的迹，由于周期性而闭合：
\begin{equation}
  P(m) = \tr\!\left[\prod_{j=0}^{N_T-1} U_4(m,j)\right].
\end{equation}

关联子 $\ev{P(m)\,P(n)^\dagger}$ 与一对静态夸克和反夸克之间的势能相关（比较 (3.61)），即在给定温度下这样一对粒子的自由能 $F_{\bar{q}q}$：
\begin{equation}
  \ev{P(m)\,P(n)^\dagger} = e^{-a N_T\,F_{\bar{q}q}(a|m-n|)}
    = e^{-F_{\bar{q}q}(r)/T}.
\end{equation}

照常，能量必须在某个距离处归一化。对于大距离，我们预期因子化
\begin{equation}
  \lim_{a|m-n|\to\infty} \ev{P(m)\,P(n)^\dagger}
    = \ev{P(m)}\,\ev{P(n)^\dagger} = \left|\ev{P}\right|^2.
\end{equation}

由于平移不变性，波利亚科夫圈的空间位置无关紧要，因此我们在 (12.7) 的右侧省略了它，而是使用空间平均
\begin{equation}
  P = \frac{1}{N^3}\sum_m P(m).
\end{equation}

我们注意到，对于随间距无限增长的静态势，例如 $\sigma > 0$ 时的禁闭势 (3.62)，$\left|\ev{P}\right|$ 必须为零。因此我们得出结论
\begin{equation}
  \ev{P} = 0 \;\Leftrightarrow\; \text{confinement},
  \qquad
  \ev{P} \ne 0 \;\Leftrightarrow\; \text{no confinement}.
\end{equation}

在低温下 QCD 是禁闭的。随着温度升高，纯规范理论在约 270~MeV 的临界温度 $T_c$ 处经历相变 [6, 7]。在那里体系退禁闭，$\ev{P}$ 获得非零值（参见图~12.1 的左侧）。

单个波利亚科夫圈 $P(m)$ 不过是一个固定在空间位置、沿时间方向延伸并因周期性而闭合的威尔逊圈。我们可以将其期望值解释为观测到单个静态电荷的概率。借助
\begin{equation}
  \left|\ev{P}\right| \sim e^{-F_q/T}
\end{equation}
我们将其与单个色荷的自由能联系起来。由 (12.9) 我们得出结论，在禁闭相中 $F_q \to \infty$，而在退禁闭相中 $F_q$ 有限。

淬火 (quenched) QCD 中的退禁闭相变在所谓的中心对称性或 $Z_3$ 对称性方面还有另一个有趣的解释。从 (3.60) 中波利亚科夫圈的定义可以明显看出，它是一个规范不变的量，不能通过规范变换使其变为 1。然而，在有限格点上，任何一个给定的规范位形与通过中心变换得到的位形具有相同的统计权重。该变换是将给定时间层 $n_4 = t_0$ 中的所有时间方向链乘以规范群 $\mathrm{SU}(3)$ 的中心群 $Z_3$ 的同一个元素 $z$：
\begin{equation}
  U_4(n,t_0) \;\to\; z\,U_4(n,t_0).
\end{equation}

$\mathrm{SU}(3)$ 的中心元素是三次单位根 $(1, e^{2\pi i/3}, e^{-2\pi i/3})$。该对称性的原因在于，规范作用量是由沿平凡闭合圈的链变量的乘积构成的。每个平凡闭合圈在一个方向上与相反方向上具有同样多的元素。中心元素与所有群元素对易，因此在闭合圈的乘积中相互抵消。例如，在 $n_4 = t_0$ 处 $(\mu,4)$ 平面内的基本方格按如下方式变换
\begin{equation}
\begin{split}
  &\tr\Big[U_\mu(n,t_0)\,U_4(n+\hat{\mu},t_0)\,U_\mu(n,t_0+1)^\dagger\,
    U_4(n,t_0)^\dagger\Big] \\
  &\quad\longrightarrow\;
  \tr\Big[U_\mu(n,t_0)\,z\,U_4(n+\hat{\mu},t_0)\,U_\mu(n,t_0+1)^\dagger\,
    U_4(n,t_0)^\dagger\,z^\dagger\Big] \\
  &\quad= \tr\Big[z\,z^\dagger\,U_\mu(n,t_0)\,U_4(n+\hat{\mu},t_0)\,
    U_\mu(n,t_0+1)^\dagger\,U_4(n,t_0)^\dagger\Big].
\end{split}
\end{equation}
由于 $zz^\dagger = 1$，规范作用量不变。空间平面内或 $n_4 \ne t_0$ 处的基本方格不受影响。这一整体对称性是第~3.4 节中强耦合极限下静态电荷禁闭证明的核心。

然而，波利亚科夫圈并非以拓扑平凡的方式闭合，而是环绕紧致时间方向一次。因此它在变换 (12.11) 下并不不变。由于我们只将一个时间层中的时间方向链乘以中心元素 $z$，波利亚科夫圈变换为
\begin{equation}
  P \;\to\; z\,P.
\end{equation}

由于作用量的对称性，波利亚科夫圈的期望值可以写为
\begin{equation}
  \ev{P} = \frac{1}{3}\left(P + zP + z^2P\right)
    = \frac{1}{3}\left(1 + e^{i2\pi/3} + e^{-i2\pi/3}\right)\ev{P} = 0 .
\end{equation}
右侧为零，因为对中心元素的求和为零。当中心对称性被破缺时，分析 (12.14) 失效。特别是，当温度超过临界值时，该对称性被自发破缺。因此，淬火 QCD 的有限温度相变标志着中心对称性的自发破缺 [8]。波利亚科夫圈（也称为热威尔逊线）是一个序参量，它区分禁闭相（其中不能找到自由电荷）与退禁闭相（其中单个电荷被屏蔽并可以被观测到）。

将多个这样的电荷组合起来，只要这些电荷结合成色单态，就会产生非零的期望值。这允许定义，例如，静态电荷之间的静态势。在具有两个或多个静态电荷的情形下，磁通管（胶子能量的集中）的演化也已被研究 [9--13]。

中心对称性类似于一个三态自旋模型的性质，该模型的自旋变量 $\in Z_3$ 位于三维空间格点体积的格点上 [14, 15]。这样的模型在某个最近邻耦合值处经历相变——但仅对无限体积而言。在有限体积中，对称性在严格意义上不能破缺。然而，在那些在无限体积情形下预期会发生对称性破缺的耦合值处进行的计算机模拟，可能在三个扇区之一中停留很长一段时间，然后才会隧穿到另一个扇区。

序参量的定义和测量是一个微妙的问题。在自旋系统中，人们引入一个外场，让体积趋于无穷，然后再让外场趋于零。这一极限顺序很重要，因为对于任何有限体积，自发磁化强度都为零。对于规范理论，我们也有类似的情况，阻碍了在有限体积中正确确定 $\ev{P}$。从技术上讲，更大的格点会导致在一个扇区中停留更长的时间，因此人们可以预期这个“准序参量”能得到较好的近似确定。然而，由于隧穿引起的系统误差无法控制。在更新过程中，甚至可以通过定期将一个时间层中的所有规范链乘以一个非平凡的中心元素来显式地强制中心对称性。这不会改变规范作用量。于是很明显，对于足够长的运行，$\ev{P}$ 必然为零。因此，人们常常转而绘制和分析 $\ev{|P|}$，后者在无限体积极限下与 $\left|\ev{P}\right|$ 一致。

为了在有限温度 $1/(aN_T)$ 下模拟，人们要么改变 $N_T$，要么改变规范耦合从而改变格点间距 $a$。改变 $a$ 允许对 $T$ 进行连续扫描。在图~12.1 中，我们展示了在大小为 $16^3\times 4$ 的格点上对纯规范理论进行模拟的结果，所选取的耦合值范围使得 $T$ 接近从禁闭到退禁闭的转变点。左侧图中我们展示了 $\ev{|P|}$ 作为 $T$ 的函数。右侧是单个位形在固定的 $T > T_c$ 值下的 $P$ 散点图。在模拟中，中心对称性也被更新，$P$ 的值集中在中心元素的方向上。

\begin{figure}[htbp]
  \centering
  \includegraphics[width=0.9\textwidth]{fig121.png}
  \caption*{图 12.1. \emph{左图：期望值 $\ev{|P|}$ 作为温度 $T$ 的函数。在计算 $T = 1/(N_T a) = 1/(4a)$ 时，格点间距是通过静态势确定的，如第~3.5 节所讨论，参见 (3.83)（其中 $r_0 = 0.5$~fm）。结果来自在格点大小 $16^3\times 4$ 上的模拟，对每个规范耦合（即温度）值平均了 100 个测量的位形；为了得到可靠的结果，应当使用更好的统计量并外推到无限体积。右图：1500 个独立规范位形的波利亚科夫圈结果在复平面上的散点图。在破缺相中，$P$ 的值集中在三个中心相位附近。数据来自 $\ev{|P|} = 0.53$ 的一次运行。}}
  \label{fig:121}
\end{figure}

结果通常会受到空间体积有限性和标度修正的影响。比较不同的参数设置可以同时研究这两者。

淬火理论的蒙特卡罗计算很早就已经完成 [16--18]，此后又得到了相当大的改进 [19, 20]。仔细的研究表明，该相变是弱一级相变。这与 $Z_3$ 自旋模型的类比一致 [14]。

在确定转变温度时，格点间距会进入，并且必须固定标度，如第~3.5 节所讨论。目前公认的淬火情形数值为 $T_c \approx 270$~MeV。来自不同改进规范作用量的结果外推到连续极限后，在百分之几的范围内一致。

在相变研究中，人们采用了统计物理的各种标准工具。一个信息量丰富的量是波利亚科夫圈极化率
\begin{equation}
  \chi_P = N^3\left(\ev{P^2} - \ev{P}^2\right),
\end{equation}
它在相变处出现峰值。该峰值的大小和位置遵循有限尺寸标度律 [21]，并通过有限体积测量给出无限体积值的信息。

有限温度下典型的计算机模拟应使空间范围远大于时间范围，以减少有限体积效应。因此 $N_T$ 相对较小，这限制了不同马祖巴频率的数目。这也使确定粒子传播子等可观测量在时间方向的指数衰减性质的尝试变得复杂，因为在大的 $T$（即小的 $N_T$）下，它们只能在很少的几个时间步内被追踪。出于这个原因，人们还在空间和时间方向分别引入了不同的格点间距 $a$ 和 $a_T$，且 $a > a_T$。这可以通过选择不同的规范耦合分别乘在时间方向基本方格和纯空间基本方格上来实现。保持比值 $\xi = a/a_T$ 较大，可以在温度恒定的同时拥有大的 $N_T$。这种在 [22, 23] 中引入的方法在 [24] 中有更详细的讨论。

\subsection{加入动力学费米子}

动力学费米子显式地破缺 $Z_3$ 对称性，这也影响波利亚科夫圈。一个简单的理解方式来自第~5.3.2 节中引入的费米子行列式的跳跃展开。(5.65) 中的指数可以被视为有效的费米子作用量。它是对闭合圈的求和，就像规范作用量本身一样。然而，这些现在都是格点上所有可能的圈，包括那些围绕紧致时间方向非平凡绕转的圈。特别是，这个求和还将包括波利亚科夫圈，因此包括与
\begin{equation}
  \kappa^{N_T}\left(P + P^*\right).
\end{equation}
成比例的项。

由于变换性质 (12.13)，有效作用量中这样的项在一个时间层中所有规范链整体乘以群的中心元素时不再不变，因此费米子破缺中心对称性。在完整 QCD 中，人们发现波利亚科夫圈趋于实的正值。在某种意义上，可以将费米子的作用看作类似于自旋模型中外部破缺对称性的磁场。

另一个更直观的图像是，动力学产生的真空圈屏蔽单个电荷，类似于可极化介质中电荷的德拜屏蔽。由关联函数 (12.6) 计算出的势能会在某个距离处饱和，并可以用屏蔽库仑势 $\propto \exp(-\mu|r|)/|r|$ 来近似。最终，“势能”这个概念失去意义，因为从真空中产生一个夸克-反夸克对来饱和试探电荷在能量上更为有利：“弦断裂”。波利亚科夫圈和威尔逊圈不再是严格的序参量。现在禁闭必须以不同的方式来定义，即通过体系的渐近态——即色中性的强子——的谱来定义。基于这些考虑，在 [25, 26] 中给出了序参量的另一种定义。

在我们的理论世界中，我们可以连续地改变夸克质量参数 $m$。让我们首先讨论具有相等质量 $m_{u,d}$ 的两味情形（参见图~12.2）。如果我们让质量趋于无穷，夸克退耦，我们就回到具有一级温度相变的淬火情形。减小夸克质量，这一相变仍然存在，但潜热减小，相变变弱。最终，$(g,m)$ 平面中的一级相变线以二级相变点结束。

\begin{figure}[htbp]
  \centering
  \includegraphics[width=0.7\textwidth]{fig122.png}
  \caption*{图 12.2. \emph{该图展示了 $\Nf = 2$ 子流形在 $(T,m)$ 平面中的相结构（图~12.3 中三维图的正面）。}}
  \label{fig:122}
\end{figure}

因此，与禁闭域和退禁闭域相邻的区域是解析连通的。这样的过渡并不罕见。只需想想水的三相：冰、液态和水蒸气。取决于压力和温度，可能存在相变，但也存在连续过渡型的转变。

然而，我们还没有结束。进一步减小夸克质量，我们最终到达无质量夸克的情形，此时作用量是手征不变的（如果使用合适的离散化）。正如我们在第~7 章所讨论的，完整理论在零温下自发破缺手征对称性。相应的序参量是在 (7.63) 中定义的夸克凝聚 $\Sigma$，然而它只是无质量理论的序参量。随着 $T$ 增大，凝聚表现出一个转变。它在某个临界温度 $T_{ch}$ 处变为零，并在更高温度下保持为零。因此，凝聚可以被视为指示手征温度转变的有效序参量（参见图~12.2）。

在第~7.3.5 节中，我们通过 Banks--Casher 关系 (7.77) 将手征凝聚与原点附近的狄拉克算子本征值谱密度联系起来。该公式也描述了转变之上的情况：在 $T_{ch}$ 处谱中打开一个间隙，手征凝聚变为零。

一个非常重要的观测结果是，禁闭似乎也在大约相同的温度 $T_{ch}$ 处消失，这表明禁闭与手征对称性破缺之间存在密切联系。$T_{deconf}$ 与 $T_{ch}$ 的相等近来受到了质疑 [1, 2, 27]。争议的可能原因可能与准临界转变点的定义以及模拟中使用的费米子种类和质量有关。

$T_{ch}$ 的值不容易确定，因为它需要在零夸克质量下进行模拟。此外，它还取决于所考虑的夸克种类数。各种模拟和外推得到的典型值范围在 $T_{ch} \approx 150$~MeV [27] 到 $T_{ch} \approx 190$~MeV [28] 之间。与纯规范理论相比，这个较小的值可以从以下观测理解：在淬火情形中，唯一的束缚态是相当重的胶球。因此，需要更高的温度来产生足够致密的胶球气体以实现退禁闭。不过，如上所述，还是应当保持一定的谨慎。

图~12.3 展示了我们目前在两个质量简并夸克（质量 $m_{u,d}$）和第三个夸克（质量 $m_s$）构成的空间中关于相图的认识。目前尚不确定，对于物理夸克质量（图中虚线圆所示），转变是位于过渡域还是位于一级相变域。

\begin{figure}[htbp]
  \centering
  \includegraphics[width=0.9\textwidth]{fig123.png}
  \caption*{图 12.3. \emph{$(T,m_{u,d},m_s)$ 空间中两个质量简并味（质量 $m_{u,d}$）加第三个味（$m_s$）、且夸克数密度/化学势为零 $\mu = 0$ 时的相结构示意图。斜线阴影区域表示过渡转变的位置，其中阴影较暗的区域表示一级相变的区域。在一级相变区域的边界上存在二级相变。位于 $m_{u,d} = 0$ 平面内的边界可能具有二级相变部分，并在三临界点处与一级相变域相接。该图基于 [29] 中所讨论和总结的信息。}}
  \label{fig:123}
\end{figure}

在物理质量值附近，具有两个轻夸克、加或不加一个较重奇异夸克的格点结果发现，从较低温到较高温存在一个快速而平滑的过渡，但没有找到令人信服的普适临界行为的证据。在更小的格点间距、更接近连续极限的条件下工作，并使用具有更小截断效应的改进费米子作用量，可能有助于进一步澄清这一状况 [30, 31]。

手征温度相变的级次不如纯规范理论那样确定，并且取决于夸克味数。Pisarski 和 Wilczek [32] 研究了序参量的一个有效三维模型场论，其假设是 QCD 潜在的自发破缺手征对称性。重整化群分析进而表明，对于 $\Nf = 2$ 手征相变是二级的，而对于 $\Nf = 3$ 则是弱一级的。在后一种情形下，相变同样向 $m > 0$ 延伸，并有一个二级端点。该端点与从纯规范理论出发的相变线的端点通过所谓的过渡曲线相连。这是一些可观测量（如比热）的最大值，但不是非解析点的轨迹。

与波利亚科夫圈类似，凝聚的导数，即手征极化率
\begin{equation}
  \chi_{ch} = \frac{\partial\Sigma}{\partial m}
\end{equation}
也是定位转变点的有用量，并且它也具有确定的有限尺寸标度行为。尽管对于 $0 < m < \infty$，$\Sigma$ 和 $\ev{P}$ 都不是严格的序参量，但它们的值仍然可以用来研究相变线的位置。

\subsection{退禁闭相中 QCD 的性质}

形式上，配分函数依赖于几个参数：体积、温度、夸克质量和规范耦合。人们通过自由能（$F = -\ln(Z)/\beta$）的导数来定义热力学量，如能量密度 $\varepsilon$ 或压强 $p$：
\begin{equation}
  \varepsilon \equiv -\frac{1}{V}\frac{\partial\ln Z}{\partial\beta}
    = \frac{T^2}{V}\frac{\partial\ln Z}{\partial T},
  \qquad
  p \equiv T\,\frac{\partial\ln Z}{\partial V},
  \qquad
  \Sigma_i = \frac{\partial\ln Z}{\partial m_i}.
\end{equation}
这里 $V$ 表示空间体积，照常 $\beta = 1/T$。

正如已经提到的，禁闭区和退禁闭区是解析连通的。尽管如此，我们仍然可以预期这些区域中的物质具有相当不同的性质，尽管转变可能是渐进的。对于大温度，跑动耦合的标度由 $T$ 设定，这使人进入小规范耦合的微扰区域。渐近自由进而使人们预期一个弱相互作用的夸克和胶子气体：等离子体。这就是为什么高温相通常被称为夸克-胶子等离子体相 [33]。

高温下（无相互作用的）QCD 胶子和夸克的相对论理想气体遵循斯特藩-玻尔兹曼定律，其压强和能量密度为 [34]，
\begin{equation}
  p = \frac{\pi^2}{45}\left(8 + \frac{21}{4}\Nf\right) T^4,
  \qquad
  \varepsilon = 3p \quad \text{(equation of state)}.
\end{equation}
由于夸克和胶子具有更多的自由度，这远大于人们从例如三个轻介子 $\pi^0$ 和 $\pi^\pm$ 的气体中所预期的值。尚不清楚强子以多快的速度解离并获得这样的行为。图~12.4 表明这可能只有在相当高的 $T$ 值下才会发生。

\begin{figure}[htbp]
  \centering
  \includegraphics[width=0.9\textwidth]{fig124.png}
  \caption*{图 12.4. \emph{从纯 SU(3) 规范理论和 QCD（两或三个轻夸克，或两个轻夸克加一个重夸克 (2+1)，其中 $m_s \approx T$ [35]）的模拟中确定的压强 $p$（图取自 [36]）。图中右侧的箭头表示来自斯特藩-玻尔兹曼定律 (12.19) 的相应无限温度极限。（来源：Karsch [36]。经 AIP 许可转载。）}}
  \label{fig:124}
\end{figure}

由于转变的过渡性质，粒子谱可能会逐渐从（零温下众所周知的）介子和重子谱变化到某个中间态，然后渐近行为开始起作用，夸克和胶子的等离子体定义了基态。$T_{ch}$ 之上手征对称性的恢复可以通过简并的宇称伙伴来实现，但仍存在例如介子束缚态。目前在重离子实验中可达的区域通常接近转变点，因此观测结果可能被中间的相变性质所主导和掩盖。

如前所述，由于欧氏时间范围很小（小于 1~fm），使用研究时间方向关联函数的标准方法来测定粒子谱相当困难，尽管已经有人尝试过 [37--39]。然而，人们可以分析空间关联函数，这导致所谓的屏蔽质量 [40--43]。将传播子对所有空间和时间积分给出热极化率，其中包含一个必须减去的接触项。不过，积分涉及相应量子通道中所有态的求和。追踪这些量穿过相变表明，$\pi(0^{-+})$ 和同位旋标量 $f_0(0^{++})$ 的质量在 $T_{ch}$ 之上很快变得简并，表明 SU(2) 味对称性恢复。然而，同位旋矢量 $a_0(0^{++})$ 的质量仍然与介子不同，表明 $U_A(1)$ 对称性破缺持续存在。

\section{化学势的引入}

真空中重子或夸克的净数为零，但在物质中并非如此。在极端情形下，如重离子碰撞或中子星中的超高密度物质，必须考虑非零夸克数密度 $n_q$ 的效应。在这种情况下，我们必须扩展巨正则系综的配分函数 (12.1)：
\begin{equation}
  Z(T,\mu) = \tr\!\left[e^{-(H - \mu N_q)/T}\right].
\end{equation}
我们向总作用量中添加了一项，引入夸克化学势 $\mu$，它乘以夸克数算符 $\hat{N}_q$。有时也使用重子数算符 $\hat{N}_B = \hat{N}_q/3$ 和重子化学势 $\mu_B = 3\mu$。

由于夸克数算符 $\hat{N}_q$ 取整数值，我们可以将巨正则配分函数 $Z(T,\mu)$ 按逸度变量
\begin{equation}
  z = e^{\mu/T}.
\end{equation}
的幂级数展开。该逸度展开的结果是固定夸克数 $n \in \mathbb{Z}$ 的正则配分函数 $Z_n(T)$ 的求和
\begin{equation}
  Z(T,\mu) = \sum_n \left(e^{\mu/T}\right)^n Z_n(T),
\end{equation}
其中 $n$ 的负值对应于反夸克的净剩余。在有限格点上，对夸克数 $n$ 的求和是有限的，即 $|n| \le n_{max}$。对于具有 $N^3$ 个点的空间体积，威尔逊费米子的每一味对 $n_{max}$ 的贡献为 $6N^3$。

借助巨正则配分函数 (12.20)，我们可以计算新的可观测量，例如夸克数密度的期望值
\begin{equation}
  n_q \equiv \frac{1}{V}\ev{\hat{N}_q}
    = \frac{T}{V}\frac{\partial\ln Z(T,\mu)}{\partial\mu}.
\end{equation}
物理环境的夸克数密度将影响相结构以及有限温度相变的位置和性质。

除了前面几章熟悉的规范耦合和夸克质量参数，以及上一节引入的温度 $T$ 之外，我们现在还有额外的参数 $\mu$。相图将变得越来越复杂，正如我们在第~12.2.2 节的讨论中将看到的那样。

\subsection{格点上的化学势}

我们如何在格点作用量中引入夸克数密度？欧几里得连续夸克数算符由守恒矢量流 $\bar{\psi}\gamma_\mu\psi$ 的时间分量
\begin{equation}
  \bar{\psi}(x)\,\gamma_4\,\psi(x)
\end{equation}
的空间体积分给出。在格点实现中，人们倾向于在狄拉克算符 (2.36) 中加入这样一项乘以 $\mu$，因为因子 $1/T$ 来自对有限时间方向的积分。然而，这种过于简单的尝试会遇到问题。为了看清这一点，让我们讨论自由情形。我们将证明，在那里能量密度 $\varepsilon(\mu)$ 在连续极限中发散，而这种非物理发散正是按 (12.24) 朴素引入化学势的后果。

自由能密度由下式给出
\begin{equation}
  \varepsilon(\mu) = -\frac{1}{V_3}\ev{\hat{H}}
    = -\frac{1}{V_3\,Z}\frac{\partial Z}{\partial\beta}
    = -\frac{1}{V_3}\frac{\partial\ln Z}{\partial\beta},
\end{equation}
其中 $V_3$ 是空间体积，$Z = \tr\exp(-\beta\hat{H})$。在格点上计算能量密度时，方便的做法是为时间方向引入一个不同的格点间距 $a_T$，在计算结束时可以将其设为与空间方向使用的格点间距 $a$ 相等。由于自由情形的配分函数是 $Z = \det[D]$，我们发现能量密度为（使用 $\beta = N_T a_T$）
\begin{equation}
\begin{split}
  \varepsilon(\mu) &= -\frac{1}{(aN)^3}\,
    \frac{\partial}{\partial(N_T a_T)}\ln(\det[D])\Bigg|_{\mu a_T=\mathrm{const}}\\
  &= -\frac{1}{(aN)^3\,N_T}\,
    \frac{\partial}{\partial(a_T)}\ln(\det[D])\Bigg|_{\mu a_T=\mathrm{const}}.
\end{split}
\end{equation}

为了计算 $\varepsilon(\mu)$，我们需要计算格点狄拉克算符 $D$ 的行列式。对于按朴素方式 (12.24) 引入化学势的自由费米子，它写为
\begin{equation}
  D(n|m) = \sum_{j=1}^3 \gamma_j\,
    \frac{\delta_{n+\hat{j},m} - \delta_{n-\hat{j},m}}{2a}
    + \gamma_4\,\frac{\delta_{n+\hat{4},m} - \delta_{n-\hat{4},m}}{2a_T}
    + m\,\delta_{nm} + \mu\,\gamma_4\,\delta_{nm}.
\end{equation}

为了计算行列式，我们将 $D$ 傅里叶变换到动量空间（参见附录~A.3）并得到
\begin{equation}
  D(p|q) = \frac{1}{|\Lambda|}\sum_{n,m\in\Lambda} e^{-ip\cdot n}\,
    D(n|m)\,e^{iq\cdot m} = \delta(p-q)\,\tilde{D}(p),
\end{equation}
其中动量空间中的狄拉克算符 $\tilde{D}$ 为
\begin{equation}
  \tilde{D}(p) = \frac{1}{a_T}\left[
    i\,\frac{a_T}{a}\sum_{j=1}^3\gamma_j\sin(p_j a)
    + i\gamma_4\sin(p_4 a_T) + a_T m\,\mathbb{1} + a_T\mu\,\gamma_4
  \right].
\end{equation}

对于 (12.26) 中所需的行列式的对数，我们得到（使用 (A.54)）
\begin{equation}
  \ln(\det[D]) = \ln\prod_p \det_d\tilde{D}(p)
    = \sum_p \ln\det_d\tilde{D}(p)
    = \sum_p \tr_d\!\left[\ln\tilde{D}(p)\right].
\end{equation}
我们使用附着在行列式和迹上的下标 $d$ 来表示二者仅对狄拉克空间中剩下的 $4\times 4$ 矩阵计算。利用 (12.26)、(12.27)、(12.28)、(12.29) 和 (12.30)，我们得到对 $a_T$ 的导数，
\begin{equation}
\begin{split}
  &\frac{\partial}{\partial a_T}\ln(\det[D])\Bigg|_{\mu a_T=\mathrm{const}}\\
  &\quad= \sum_p \tr_d\!\left[
    \frac{\partial}{\partial a_T}\ln\tilde{D}(p)
  \right]_{\mu a_T=\mathrm{const}} \\
  &\quad= C + \sum_p \tr_d\!\left[
    \left(a_T\tilde{D}(p)\right)^{-1}\left(
      \frac{i}{a}\sum_{j=1}^3\gamma_j\sin(p_j a) + m\,\mathbb{1}
    \right)
  \right],
\end{split}
\end{equation}
其中常数 $C$ 由 $C = -4N_T N^3/a_T$ 给出。计算狄拉克空间中其余的迹并将结果代入 (12.26)，在令 $a_T = a$ 之后我们得到
\begin{equation}
  \varepsilon(\mu) = C - \frac{4}{N^3 N_T a^4}\sum_p F(ap, am, a\mu),
\end{equation}
其中
\begin{equation}
  F(ap, am, a\mu)
  = \frac{\sum_{j=1}^3\sin^2(a p_j) + (am)^2}
         {\sum_{k=1}^3\sin^2(a p_k) + (am)^2
          + \left(\sin(a p_4) - i\,a\mu\right)^2}.
\end{equation}

在计算能量密度的最终结果时，我们需要通过减去 $\mu = 0$ 的值来归一化，即考虑 $\varepsilon(\mu) - \varepsilon(0)$。我们回到零温，并令 $N = N_T \to \infty$，将对动量 $p$ 的求和改为连续积分：
\begin{equation}
  \frac{1}{N^3 N_T a^4}\sum_p
  \;\Rightarrow\;
  \frac{1}{(2\pi)^4}\int_{-\pi}^{\pi}\dd^4 p
\end{equation}
从而得到
\begin{equation}
  \varepsilon(\mu) - \varepsilon(0) = -\frac{4}{(2\pi)^4}\int_{-\pi}^{\pi}
    \dd^4 p\;\big(F(ap,am,a\mu) - F(ap,am,0)\big),
\end{equation}
其中在 $F$ 中我们仍然保持 $a$ 有限。使用复 $p_4$ 平面中的围道积分，可以提取出 $a\to 0$ 极限下的首项。由 $p_4 \to 0$ 的贡献得到行为
\begin{equation}
  \lim_{a\to 0}\big(\varepsilon(\mu) - \varepsilon(0)\big) \propto
    \left(\frac{\mu}{a}\right)^2 .
\end{equation}

这在连续极限中是发散的，因此在作用量中简单地加上项 (12.24) 并不能给出化学势的合适离散化。

对连续情形的仔细考察（参见 [44--46]）澄清了这个问题。夸克数是 U(1) 整体对称性的守恒荷。对格点作用量确定 Noether 流，得到的流由最近邻项表示。空间积分进而产生合适形式的化学势项。因此，人们通过将 (5.51) 中的时间方向跳跃项替换为下式来实现化学势
\begin{equation}
\begin{split}
  -\frac{1}{2a}\sum_{n\in\Lambda}\Big[
  &f(a\mu)\,(1-\gamma_4)_{\alpha\beta}\,U_4(n)_{ab}\,\delta_{n+\hat{4},m}\\
  &\quad - f(a\mu)^{-1}\,(1+\gamma_4)_{\alpha\beta}\,
    U_4(n-\hat{4})^\dagger_{ab}\,\delta_{n-\hat{4},m}
  \Big],
\end{split}
\end{equation}
其中 $f(a\mu)$ 是一个尚未指定的实函数。

对于 $\mu = 0$ 应恢复原始作用量，因此我们要求 $f(0) = 1$。在 $\mu$ 的展开中，下一项应线性依赖于化学势，以便重现密度项。因此我们有 $f(a\mu) = 1 + a\mu + \mathcal{O}(a\mu)^2$。时间反射不变性（参见第~5.4 节）要求 $f(a\mu) = 1/f(-a\mu)$。满足所有这些条件的最简单选择是
\begin{equation}
  f(a\mu) = \exp(a\mu).
\end{equation}

在这种表述中，时间方向的前向传播被因子 $\exp(a\mu)$ 增强，而向后传播被 $\exp(-a\mu)$ 抑制。这引入了所需的粒子-反粒子不对称性。

在第~5.3.1 节中，我们展示了跳跃展开将费米子行列式表示为对闭合圈的求和。在这样的圈中，前向跳跃因子 $f(a\mu)$ 和反向跳跃因子 $f(a\mu)^{-1}$ 相互抵消，除非该圈非平凡地绕转紧致时间方向。化学势的总贡献则为 $f(a\mu)^{wN_T}$，其中 $w \in \mathbb{Z}$ 表示圈的绕转数。

基于这一观测，我们注意到在模拟中，除了修改时间方向的所有链项之外，也可以只在单个时间层中用因子
\begin{equation}
  f(a\mu)^{N_T} = \exp(a\mu N_T) = \exp(\mu/T),
\end{equation}
修改所有前向时间方向的跳跃项，并用其逆因子修改相应的后向项。与 (12.20) 中一样，我们发现化学势总是以 $\mu/T$ 的形式出现。

化学势的引入带来一个严重的技术缺点：对于 $a\mu \ne 0$，狄拉克算符不再 $\gamma_5$-厄米。容易看出，将跳跃项 (12.37) 从左侧和右侧乘以 $\gamma_5$ 会改变 $\gamma_4$ 的符号；厄米共轭随后交换 $U$ 和 $U^\dagger$，我们发现代替 $\gamma_5$-厄米关系 (5.76) 的是修正的方程
\begin{equation}
  \gamma_5\,D(\mu)\,\gamma_5 = D^\dagger(-\mu).
\end{equation}

然而，这不再是一个对称性变换，而是导致一个将 $f$ 替换为 $1/f^*$ 的项。我们得到
\begin{equation}
  \gamma_5\,D(f)\,\gamma_5 = D^\dagger(1/f^*)
  \quad\Rightarrow\quad
  \det[D(f)] = \left(\det[D(1/f^*)]\right)^* .
\end{equation}

这只有在 $f = 1/f^*$ 时才是不变性操作。对于实的 $f$，因此有 $f = 1/f = 1$，这意味着 $a\mu = 0$。因此，对于非零的实 $\mu$，狄拉克算符的行列式是复的，并且不能通过将费米子味数加倍来获得实的玻尔兹曼权重。非零的实 $\mu$ 产生粒子-反粒子不对称性，它破坏了行列式的实性以及重要性采样的直接应用。

然而，有趣的是注意到，对于虚化学势 $\mu = i\eta$，$\eta \in \mathbb{R}$，我们有
\begin{equation}
  f(ia\eta) = 1/f(ia\eta)^* = 1/f(-ia\eta).
\end{equation}

于是狄拉克算符以标准方式 $\gamma_5$-厄米，行列式为实 [47--49]。基于这一观测，人们可以尝试在虚化学势下工作，然后解析延拓到实值（参见第~12.3.3 节）。

另一个特殊情形是考虑所谓的同位旋化学势 $\mu_I$ [49, 50]。让我们为每一味引入一个化学势变量。在欧几里得连续拉格朗日量中，这对应于项
\begin{equation}
  \sum_f \mu_f\,\bar{\psi}_f\,\gamma_4\,\psi_f,
\end{equation}
可以按照上面的讨论将其转换到格点表述。关注两个轻味，我们可以取 $\mu_u = \mu_I$ 和 $\mu_d = -\mu_I$。拉格朗日量中的项则具有形式 $\mu_I(\bar{u}\gamma_4 u - \bar{d}\gamma_4 d)$。在这种情况下，狄拉克算符是块对角的，其中块是上夸克和下夸克的单味算符：
\begin{equation}
  \begin{pmatrix} D(\mu_I) & 0 \\ 0 & D(-\mu_I) \end{pmatrix}
  =
  \begin{pmatrix} D(\mu_I) & 0 \\
                  0 & \gamma_5\,D^\dagger(\mu_I)\,\gamma_5 \end{pmatrix}.
\end{equation}
在这一步中我们使用了 (12.40)。对于质量简并的夸克，该矩阵的行列式是实且正的，因为
\begin{equation}
  \det[D(\mu_I)]\,\det[\gamma_5\,D^\dagger(\mu_I)\,\gamma_5]
  = \det[D(\mu_I)]\,\det[D^\dagger(\mu_I)]
  = \left|\det[D(\mu_I)]\right|^2 .
\end{equation}
这个性质允许人们使用第~8 章中讨论的标准技术，在蒙特卡罗模拟中研究非零同位旋化学势。同位旋化学势并不是一个纯学术问题，因为可以想象在相对论重离子碰撞中会出现这样的情况。

对于非零化学势，也存在重叠狄拉克算符的定义 [51, 52]。核算符可以是按 (12.37) 引入化学势 $\mu$ 的通常威尔逊狄拉克算符，它不再是 $\gamma_5$-厄米的。在构造 $D_{ov}$ 时应用的符号函数（参见第~7.4 节）是对非厄米但可对角化算符的推广，即在谱表示 (7.84) 中取复本征值实部的符号。对于 $\mu \ne 0$，$D_{ov}(\mu)$ 不是 $\gamma_5$-厄米的，但继承了性质 (12.40)。它不再是正规算符，其本征值也不限于 Ginsparg--Wilson 圆。然而，本征矢量 $\psi$ 和 $\gamma_5\psi$ 的本征值仍然通过 $\lambda$ 和 $\lambda/(\lambda - 1)$ 相关联。已有证明 [53, 54] 表明，这一推广正确地再现了自由费米气的热力学性质。然而，在存在非零化学势时，手征对称性在格点上的精确体现仍在讨论之中。

\subsection{$(T,\mu)$ 空间中的 QCD 相图}

关于相结构细节的情况，如相变的精确位置和临界指数值，仍有待确定。来自不同类型狄拉克算符的格点结果常常不一致，数值和统计量可能不足。因此，这里讨论的相结构只能代表目前的最低共识 [55, 56]。

图~12.5 展示了两个无质量夸克和两个接近物理质量的轻夸克的情形。

\begin{figure}[htbp]
  \centering
  \includegraphics[width=0.9\textwidth]{fig125.png}
  \caption*{图 12.5. \emph{$(T,\mu)$ 平面中 (左图) 两个无质量夸克和 (右图) 两个质量简并的轻夸克的推测相图。迄今为止的格点计算仅对化学势较小的值有用，试图定位一级相变线的三临界点 (左图) 或二级端点 (右图)。}}
  \label{fig:125}
\end{figure}

$\mu = 0$ 处的温度相变已在第~12.1.2 节中讨论过：对于 $\Nf = 2$，它要么是二级的，要么是弱一级的。随着夸克质量增大，相变变弱，最终变为过渡。通过引入一个轻的第三个动力学夸克，相变变强，从相图中部出发的一级相变线向 $\mu = 0$ 区域延伸。如果第三个夸克的质量增大，则又回到图~12.5 右侧图的情形。

在零温和非零实化学势下没有格点结果。高密度下 QCD 的直接微扰处理也失效 [55]。那里的大部分信息来自受 NJL 启发的计算 [57]、具有各种相互作用模型的 Dyson--Schwinger 间隙方程，如瞬子液体模型 [58] 和随机矩阵模型 [59, 60]。在禁闭相中，首先存在核物质，其（对两个无质量味而言）有一条小的相变线延伸入相中，将类共振气体结构与类液体结构分开。模型计算进而表明，在 $\mu_B$ 高于 1~GeV 的核物质区域与更高密度下的（两味）色超导相 (2SC) 之间存在一级相变，在 2SC 相中夸克在味单态色反三重态道中配对 [57, 58]。加入一个质量接近其物理值的奇异夸克后，相变周围的区域可能更加复杂，例如允许出现一个结晶色超导相的中间口袋。然而，所有计算都一致认为，对于较大的 $\mu$，人们最终进入具有非零双夸克凝聚的色-味锁定 (CFL) 相 [61]。

迄今为止，格点结果仅限于接近零化学势的窄条带，通常 $\mu/T < 1$。所有方法都基于从 $\mu = 0$ 或虚 $\mu$ 到非零 $\mu$ 的解析延拓。对重离子碰撞特别重要的是，物理情形是接近过渡还是相变线，以及在各种实验中最终处于哪一相。

\section{化学势：蒙特卡罗技术}

寻找非零化学势的最优模拟策略远未结束。因此，我们不试图涵盖所有可用的结果，而只讨论基本概念。到目前为止，所有方法都是将实行列式测量所得的结果解析延拓到人们感兴趣的实际参数值。这样的外推有两类：
\begin{itemize}
  \item 使用 $\mu = 0$ 处的测量结果，借助重整加权或 $\mu/T$ 的幂级数（矩展开）进行外推。
  \item 使用纯虚 $\mu$（负 $\mu^2$）下确定的结果，通过幂级数和 Pad\'e 有理展开，或通过傅里叶变换重构逸度展开系数，解析延拓到实的 $\mu$（正 $\mu^2$）。
\end{itemize}

淬火模拟 [62--64] 最初导致了一些令人困惑的结果。人们发现 $\mu$ 的临界值随介子质量减小：$\mu_c \propto m_\pi/2$，这在手征极限下将变为零。另一方面，在零温下预期转变在 $\mu_c \approx m_p/3$ 附近，因为质子是拥有非零重子数的最轻态。Stephanov [65] 指出，对于 $\mu = 0$，淬火理论并不是 QCD 的 $\Nf \to 0$ 极限，而是具有 $\Nf$ 个夸克和 $\Nf$ 个共轭夸克的理论的极限。这意味着在模拟中夸克行列式的相位很重要，必须使用动力学费米子才能得到 $\mu_c \approx m_p/3$。因此，该背景下的大多数模拟都使用动力学费米子，主要是 staggered（交错）类型的费米子。

大温度 $T \gg T_c$ 对应于非常小的时间范围，人们预期体系近似趋向一个三维规范理论。时间圈则扮演伴随希格斯场的角色 [66, 67]。这种维度约化开辟了在有效理论中进行包括化学势在内的蒙特卡罗研究的道路 [68--71]。我们将讨论限制在完整四维体系的方法上。

\subsection{重整加权}

重整加权是统计自旋系统蒙特卡罗方法中的标准工具，已成功用于改进不同耦合下蒙特卡罗结果之间的插值，以及从实耦合到复耦合的解析延拓。它在确定复耦合下配分函数零点（Lee--Yang 零点和 Fisher 零点）的位置方面尤其有用。

该策略可以被理解为一个更一般的概念，我们这里用一个简单的玻色子系统简要讨论。配分函数被改写为对体系总能量可能值的求和：
\begin{equation}
\begin{split}
  Z(\kappa) \equiv \int \mathcal{D}[\Phi]\,e^{-\kappa S[\Phi]}
    = \int\dd E\,\rho(E)\,e^{-\kappa E},
  \qquad \text{with} \\
  \rho(E) = \int \mathcal{D}[\Phi]\,\delta(E - S[\Phi]).
\end{split}
\end{equation}

这样，能量分布密度 $\rho(E)$ 可以通过组合来自不同耦合值的蒙特卡罗结果以高精度获得——即所谓的多直方图技术 [72, 73]。一旦态密度 $\rho(E)$ 已知，就可以对任何耦合值 $\kappa$ 计算 $Z(\kappa)$ 及其导数。然而，这只有在能量分布被精确到任意精度时才正确。在实际模拟中并非如此，因此技术上的关键点是，从某些耦合处的蒙特卡罗运行生成的分布与想要计算 $Z$ 的其它耦合值处的分布重叠程度。

本着将行列式用作权重因子的精神，可以使用 $\det[D(\mu)]$ 与 $\det[D(0)]$ 的偏差作为重整加权因子。由
\begin{equation}
\begin{split}
  &\left\langle
    \frac{e^{-S_G[\beta',U]}\,\det[D(\mu',U)]}
         {e^{-S_G[\beta,U]}\,\det[D(\mu,U)]}
  \right\rangle_{\beta,\mu}\\
  &\quad= \frac{1}{Z(\beta,\mu)}\int\mathcal{D}[U]\,
     \frac{e^{-S_G[\beta',U]}\,\det[D(\mu',U)]}
          {e^{-S_G[\beta,U]}\,\det[D(\mu,U)]}\,
     e^{-S_G[\beta,U]}\,\det[D(\mu,U)] \\
  &\quad= \frac{1}{Z(\beta,\mu)}\int\mathcal{D}[U]\,
     e^{-S_G[\beta',U]}\,\det[D(\mu',U)]
  = Z(\beta',\mu')/Z(\beta,\mu),
\end{split}
\end{equation}
我们发现，至少在理论上，可以用在 $\beta$ 和 $\mu$ 处评估的结果来确定 $Z(\beta',\mu')$。因此我们可以得到
\begin{equation}
  Z(\beta',\mu') = Z(\beta,\mu)\,
  \left\langle
    \frac{e^{-S_G[\beta',U]}\,\det[D(\mu',U)]}
         {e^{-S_G[\beta,U]}\,\det[D(\mu,U)]}
  \right\rangle_{\beta,\mu}.
\end{equation}
的值。

（在这些方程中，$\beta$ 表示逆规范耦合，而不是逆温度。）由于在蒙特卡罗抽样中，位形是以由 $\beta$ 和 $\mu$ 值决定的概率获得的，结果关键取决于人们想要外推多远。在研究有限温度相变时，人们在 $(\beta = \beta_c, \mu = 0)$ 附近抽样，其中 $\beta_c$ 表示使得 $aN_T = 1/T_c$ 的规范耦合值。

技术问题在于在 $\beta,\mu$ 处抽样的系综与目标耦合 $\beta',\mu'$ 处实际系综的重叠，以及由此产生的统计可靠性。总自由能的差异随格点体积增大而增大，使得更大格点的外推更加困难。随着化学势增大，情况变得更糟，预测能力更差。

最初的尝试保持规范耦合固定 [62, 64]，从小格点上零化学势的模拟外推到非零值。

Fodor 和 Katz [74, 75] 通过在两个变量（规范耦合和化学势）上重整加权来外推，在转变点 $(\beta = \beta_c, \mu = 0)$ 处抽样，并沿准临界线推进。该策略似乎改进了近似，并已被应用于小格点。分析配分函数的复（Lee--Yang）零点可以定位预期一级相变线的端点。在 2+1-staggered 夸克接近其物理质量值的模拟中，该点被估计位于 $T \approx 162(2)$~MeV 处的 $\mu_c^E \approx 120(13)$~MeV [76]，对应于 $\mu_c^E/T \approx 0.75$。

对于虚化学势，可以用标准方法进行模拟，因此可以对重整加权方法的收敛性质进行系统检验。关于这种比较的批判性评估，参见 [77--79]。

\subsection{级数展开}

化学势总是以组合 $\mu/T$ 出现。在另一种外推方法中，人们因此将物理可观测量写成 $\mu/T$ 的幂级数。规范位形系综的时间反演不变性（比较第~5.4 节）提供了对称性
\begin{equation}
  Z(\mu/T) = Z(-\mu/T).
\end{equation}

因此 $Z$ 和那些在时间反演下对称的可观测量是 $\mu/T$ 的偶函数，展开将按偶次幂 $(\mu/T)^2$ 进行。均匀系统中的压强密度 (12.18) 为 $p = (T/V)\ln Z$。相应的 Taylor--McLaurin 级数在 $\mu$ 下是偶的，写为
\begin{equation}
  p = T^4\sum_{n=0}^{\infty} c_{2n}(T)\left(\frac{\mu}{T}\right)^{2n},
\end{equation}
其中系数是自由能 $F = -\ln(Z)/\beta$ 的导数，因此是广义夸克数极化率，在 $\mu$ 为零处取值 [80--83]。

(12.50) 中二次项的系数是 $\mu = 0$ 处的夸克数极化率。它是夸克数密度 (12.23) 的导数，也是费米子行列式对化学势的二阶导数，
\begin{equation}
  \chi_q = \frac{\partial n_q}{\partial\mu}
    = \frac{T}{V}\frac{\partial^2\ln Z}{\partial\mu^2}.
\end{equation}
这样的导数涉及如下迹
\begin{equation}
  \frac{\partial\ln\det[D]}{\partial\mu}
  = \tr\!\left[D^{-1}\frac{\partial D}{\partial\mu}\right]
\end{equation}
以及更高阶的导数（参见 [83]）。它们的计算方式与 (8.44) 和 (8.47) 中费米子力背景下讨论的类似。迹用第~8.4 节讨论的随机估计器来评估。(12.50) 中的系数是在 $\mu = 0$ 的模拟中确定的。

在热力学极限下，相变将与某些可观测量的奇异性相关。展开的系数可以用来估计最近奇异点的位置。夸克数极化率在相变处表现出显著的峰 [84]。有限尺寸标度考虑进而可能给出相变或过渡类型的提示。几个小组（特别是 [82--86]）以这种方式分析了夸克数极化率和其他展开系数。

\subsection{虚 $\mu$}

在 (12.42) 中我们注意到，对于虚化学势 $\mu = i\eta$，$\eta \in \mathbb{R}$，狄拉克算符是 $\gamma_5$-厄米的，行列式为实。于是可以像无化学势的情形那样应用标准的重要性采样技术。

化学势通过因子 $\exp(\pm i\eta/T)$ 进入配分函数，这可以从 (12.39) 看出。因此热力学函数在 $\eta/T$ 下是 $2\pi$ 周期的，我们可以将 $\eta/T$ 限制在区间 $[0, 2\pi)$ 内。

让我们讨论复 $z$ 平面中的解析结构 [87]，其中
\begin{equation}
  z \equiv (\mu/T)^2 = -(\eta/T)^2 .
\end{equation}

动力学费米子破缺纯规范理论的 $Z_3$ 对称性（见第~12.1.2 节）。然而，$Z_3$ 变换对费米子场的影响可以通过复化学势的平移来补偿，从而导致一个新的对称性。这意味着化学势虚部具有周期性，
\begin{equation}
  Z(\eta/T) = Z(\eta/T + 2\pi/3) .
\end{equation}

因此区间 $\eta/T \in [0, 2\pi)$ 被分成三个等价扇区，它们仅由波利亚科夫圈的相位来区分。对于虚化学势，在扇区边界处存在一级的 $Z_3$ 相变，以及在某个温度之上存在过渡线 [88, 89]。这些奇异性将级数的收敛域限制在约 $|\mu/T| \le 1$。

人们可以直接在这样的负 $z$ 值下模拟该理论，即虚化学势 $\mu = i\eta$，并尝试将结果外推到正的 $z$ [88--92]。解析延拓可以通过多种方式进行。最简单的方法是在 $z$ 中的幂级数，它在一个圆形区域 $|z - z_0| < R$ 内收敛，其中 $R$ 是到最近奇异点的距离。解析函数理论中的一种标准方法是在前一级数收敛域内合适选取的点周围进行一系列展开。其他方法包括最优映射或 Pad\'e 展开。后者是一种系统的方法，用具有相同展开系数的有理函数替换幂级数。

一种直接的方法是将若干 $z$ 值处的模拟结果拟合到如 (12.50) 那样的幂级数。Taylor 展开与 [81, 83] 中针对压强所示具有相同的收敛性质。结果表明，对于所研究的几个可观测量，如手征凝聚和屏蔽质量 [69]，以及过渡的位置（准临界温度）[79, 88, 89, 92]，收敛是快速的。Pad\'e 近似可能允许将外推范围扩展到幂级数收敛圆之外 [93]。

另一种对复 $\mu$ 值的扩展引入了变量 $\cosh(\mu a)$ 和 $\sinh(\mu a)$ [94, 95]，随后解析外推到实化学势。过渡位置的结果与上述工作一致。

\subsection{正则配分函数}

对于虚化学势，巨正则配分函数的逸度展开 (12.22) 是一个傅里叶求和 [47, 48]，并且关于 $\eta/T$ 以 $2\pi$ 为周期，
\begin{equation}
  Z(T, i\eta) = \sum_n e^{i n\eta/T}\,Z_n(T).
\end{equation}

这个求和可以求逆，从而得到固定夸克数 $n$ 下正则配分函数 $Z_n(T)$ 的如下表达式：
\begin{equation}
  Z_n(T) = \frac{1}{2\pi}\int_{-\pi}^{\pi}\dd(\eta/T)\;
    e^{-i n\eta/T}\,Z(T, i\eta).
\end{equation}
因此方程 (12.56) 是逸度展开 (12.22) 系数 $Z_n$ 的精确表达式。

在虚化学势下的模拟可以用来确定 $Z(T,i\eta)/Z(T,0)$ 的值，因此原则上可以数值计算 $Z_n$ 的傅里叶积分 (12.56)。一个问题是傅里叶求和的振荡随夸克数 $n$ 增大而加剧，因此在较大的 $|n|$ 值处，为了对傅里叶积分进行可靠的数值估计，需要在区间 $[-\pi,\pi]$ 内取许多 $\eta/T$ 的中间点。可以结合多直方图采样和其他精细的方法来确定 $Z_n$，然后通过 (12.22) 得到 $Z(T,\mu)$ [78, 96--101]。最近探索的一种替代方法 [102, 103] 是用微扰技术来计算 $Z_n$。

研究正则配分函数特别适合探索低温下的少数核子系统，并且原则上允许研究核物质的体性质和核相互作用。亚稳态区域被很好地展现出来，Maxwell 构造将这个信号与巨正则结果联系起来。这种方法得到的结果与前面讨论的其他方法的结果之间良好的一致性已被证明 [98]。迄今为止的结果都在小格点上。不过，固定重子密度要求随着体积增大而增大 $n$。这对于模拟所需的统计量来说可能是一个问题。

% [unrecovered: the reference list in extract/ch12.txt is truncated after
% entry 52; entries 53-103, which are cited in the text, were not extracted]

\section*{参考文献}

\begin{enumerate}[label={[\arabic*]}]
  \item Z.~Fodor: PoS LATTICE2007, 011 (2007) 301, 308.
  \item F.~Karsch: PoS LATTICE2007, 015 (2007) 301, 308.
  \item O.~Philipsen: in \emph{Conceptual and Numerical Challenges in Femto-
        and Peta-Scale Physics}, edited by C.~Gattringer et al., The European
        Physical Journal ST 152, p.~29 (Springer, Berlin, Heidelberg, New York
        2007) 301.
  \item C.~DeTar: PoS LATTICE2008, 001 (2008) 301.
  \item S.~Ejiri: PoS LATTICE2008, 002 (2008) 301.
  \item A.~M.~Polyakov: Phys.\ Lett.\ B 72, 477 (1978) 303.
  \item L.~Susskind: Phys.\ Rev.\ D 20, 2610 (1979) 303.
  \item G.~'t Hooft: Nucl.\ Phys.\ B 153, 141 (1979) 305.
  \item C.~B.~Lang and M.~Wiltgen: Phys.\ Lett.\ B 131, 153 (1983) 305.
  \item G.~Bali: Phys.\ Rep.\ 343, 1 (2001) 305.
  \item H.~Ichie, V.~Bornyakov, T.~Streuer, and G.~Schierholz: Nucl.\ Phys.\ A
        721, 899 (2003) 305.
  \item F.~Bissey et al.: Phys.\ Rev.\ D 76, 114512 (2007) 305.
  \item Y.~Peng and R.~W.~Haymaker: Phys.\ Rev.\ D 47, 5104 (1993) 305.
  \item B.~Svetitsky and L.~G.~Yaffe: Nucl.\ Phys.\ B 210, 423 (1982)
        305, 306.
  \item B.~Svetitsky and L.~G.~Yaffe: Phys.\ Rev.\ D 26, 963 (1982) 305.
  \item L.~D.~McLerran and B.~Svetitsky: Phys.\ Lett.\ B 98, 195 (1981) 306.
  \item J.~Kuti, J.~Polonyi, and K.~Szlachanyi: Phys.\ Lett.\ B 98, 199 (1981)
        306.
  \item L.~D.~McLerran and B.~Svetitsky: Phys.\ Rev.\ D 24, 450 (1981) 306.
  \item E.~Laermann and O.~Philipsen: Ann.\ Rev.\ Nucl.\ Part.\ Sci.\ 53, 163
        (2003) 306.
  \item F.~Karsch and E.~Laermann: in \emph{Quark Gluon Plasma}, edited by\\
        R.~C.~Hwa, p.~1 (World Scientific, Singapore 2003) 306.
  \item J.~L.~Cardy: \emph{Finite-Size Scaling} (North-Holland, Amsterdam
        1988) 306.
  \item J.~Engels, F.~Karsch, I.~Montvay, and H.~Satz: Phys.\ Lett.\ B 101, 89
        (1981) 307.
  \item J.~Engels, F.~Karsch, I.~Montvay, and H.~Satz: Nucl.\ Phys.\ B 205
        [FS5], 545 (1982) 307.
  \item I.~Montvay and G.~M\"unster: \emph{Quantum Fields on a Lattice}
        (Cambridge University Press, Cambridge, New York 1994) 307.
  \item K.~Fredenhagen and M.~Marcu: Commun.\ Math.\ Phys.\ 92, 81 (1983) 307.
  \item J.~Bricmont and J.~Fr\"ohlich: Phys.\ Lett.\ B 122, 73 (1983) 307.
  \item Y.~Aoki, Z.~Fodor, S.~D.~Katz, and K.~K.~Szabo: Phys.\ Lett.\ B 643, 46
        (2006) 308.
  \item M.~Cheng et al.: Phys.\ Rev.\ D 74, 054507 (2006) 308.
  \item F.~Karsch: Lect.\ Notes Phys.\ 583, 209 (2002) 309.
  \item Y.~Aoki, Z.~Fodor, S.~D.~Katz, and K.~K.~Szabo: JHEP 0601, 089 (2006)
        309.
  \item C.~Bernard et al.: PoS LAT2005, 156 (2005) 309.
  \item R.~D.~Pisarski and F.~Wilczek: Phys.\ Rev.\ D 29, 338 (1984) 309.
  \item E.~Shuryak: Phys.\ Rep.\ 61, 71 (1980) 310.
  \item R.~P.~Feynman: \emph{Statistical Mechanics} (Westview Press, Boulder
        1998) 310.
  \item F.~Karsch, E.~Laermann, and A.~Peikert: Phys.\ Lett.\ B 478, 447 (2000)
        311.
  \item F.~Karsch: in \emph{AIP Conf.\ Proc.\ (PANIC05)}, Vol.~842, p.~20
        (Am.\ Inst.\ Phys., Melville, New York 2006) 311.
  \item P.~de Forcrand et al.: Phys.\ Rev.\ D 63, 054501 (2001) 311.
  \item I.~Wetzorke et al.: Nucl.\ Phys.\ B (Proc.\ Suppl.) 106, 510 (2002)
        311.
  \item M.~Asakawa, T.~Hatsuda, and Y.~Nakahara: Nucl.\ Phys.\ A 715, 863
        (2003) 311.
  \item J.~B.~Kogut, J.~F.~Lagae, and D.~K.~Sinclair: Phys.\ Rev.\ D 58,
        054504 (1998) 311.
  \item I.~Pushkina et al.: Phys.\ Lett.\ B 609, 265 (2005) 311.
  \item S.~Wissel et al.: PoS LAT2005, 164 (2005) 311.
  \item R.~V.~Gavai, S.~Gupta, and R.~Lacaze: PoS LAT2006, 135 (2006) 311.
  \item P.~Hasenfratz and F.~Karsch: Phys.\ Lett.\ B 125, 308 (1983) 315.
  \item J.~Kogut, H.~Matsuoka, M.~Stone, and H.~W.~Wyld: Nucl.\ Phys.\ B 225
        [FS9], 93 (1983) 315.
  \item N.~Bilic and R.~V.~Gavai: Z.\ Phys.\ C 23, 77 (1984) 315.
  \item E.~Dagotto, A.~Moreo, R.~L.~Sugar, and D.~Toussaint: Phys.\ Rev.\ B 41,
        811 (1990) 316, 322.
  \item A.~Hasenfratz and D.~Toussaint: Nucl.\ Phys.\ B 371, 539 (1992)
        316, 322.
  \item M.~G.~Alford, A.~Kapustin, and F.~Wilczek: Phys.\ Rev.\ D 59, 054502
        (1999) 316.
  \item D.~T.~Son and M.~A.~Stephanov: Phys.\ Rev.\ Lett.\ 86, 592 (2001) 316.
  \item J.~Bloch and T.~Wettig: Phys.\ Rev.\ Lett.\ 97, 012003 (2006) 317.
  \item J.~Bloch and T.~Wettig: Phys.\ Rev.\ D 76, 114511 (2007) 317.
\end{enumerate}
